\scanpage{141}
% archive foliation: 140

La catégorie des groupes finis \inserted{comm. \(G\)} sur \(K\) tels que pour
\(\forall l\), \(G(l)\) soit non ramifié rel. à \(S[1/l]\), est une sous
catégorie épaisse de celle des motifs sur \(S\), la catégorie
quotient est équivalente à celle des isomotifs lisses sur \(S\).

La catégorie des motifs est une \(\otimes\)-catégorie avec
\uline{Hom}. Les foncteurs Motifs \(\longrightarrow\) isomotifs, motifs \(\longrightarrow\)
faisceaux \(l\)-adiques, sont des \(\otimes\)-foncteurs, compatibles aux
\uline{Hom}.

\unsure{Un|On} \deleted{définit} hom de motifs est une isogénie, s'il induit
un isom sur les isomotifs sous-jacents. On définit
la dimension (ou le poids), le niveau etc en termes
des isomotifs.

Tout schéma propre et lisse \(X\) sur \(S\) définit des
motifs \inserted{sur \(S\)} notés \(R^if_*(\mathbb Z_X)\), s'il n'y a pas de confusion à
craindre. De \(=\) pour les \unsure{sites} étales relatifs lisses et propres.

\uline{Relations avec les schémas en groupes commutatifs propres et lisses.}

On définit un foncteur \uline{contravariant} de la catégorie de
ces objets, soit \(A\), dans la catégorie des motifs sur \(S\),
en posant
\[
\deleted{\unsure{\varphi|\Phi}(A)=R^1f_*(\mathbb Z_A)=
(\quad R^1\unsure{f_*|f^*}(\quad}
\]
\[
\unsure{\varphi|\Phi}(A)=(M,E_l,u)
\]
où
\[
\left\{
\begin{aligned}
M&=R^1(f_{A^0})_*(\mathbb Q_{A^0})\\
E_l&=\underline{\operatorname{Hom}}({}_{l^\infty}A,\mathbb Q_l/\mathbb Z_l)\\
u_l&\ \text{provenant de l'hom évident}\\
&\qquad E_l\longrightarrow R^1f_{A^0*}(\mathbb Z_{lA^0})
\end{aligned}
\right.
\]
\margin{extension de \(\underline{\operatorname{Hom}}(A/A^0,\mathbb Q_l/\mathbb Z_l)\)
par \(R^1f_{A^0*}(\mathbb Z_{lA^0})\).}

\uline{NB} Il suffirait que \(A\)
\deleted{\(\unsure{/A^0|A^0}\) \illegible{2}}
\inserted{\unsure{soit} extension d'un schéma}
\unsure{gr. fini|en gr.} \(A/A^0\) par un schéma abélien \(A^0\), avec
\(A/A^0(l)\) extension par \(0\) d'un groupe fini étale sur \(S_{\unsure{l}}\).
